\section{Grenzen und Limitationen}

Das entwickelte ERP-KMU-Fit-Modell dient dazu, zentrale Voraussetzungen für die wirksame Nutzung von ERP-Systemen in 
KMU strukturiert sichtbar zu machen. Es stellt jedoch keine allgemeingültige Erfolgsformel dar. Vielmehr ist es als 
konzeptionelles Orientierungsinstrument zu verstehen, das Unternehmen bei der Einschätzung ihrer Ausgangslage und 
möglicher Handlungsbedarfe unterstützen kann.

Eine erste Grenze liegt in der notwendigen Vereinfachung betrieblicher Realität. Die fünf Dimensionen des Modells 
ermöglichen zwar eine strukturierte Betrachtung von Prozessen, Daten, Ressourcen, Kompetenzen und Strategie. In der 
Praxis sind diese Bereiche jedoch nicht unabhängig voneinander. Beispielsweise kann eine geringe Ressourcenverfügbarkeit 
die Qualität von Schulungen beeinträchtigen und dadurch zugleich den Kompetenz-Fit schwächen. Ebenso können unklare 
Prozesse zu Problemen bei der Datenpflege führen. Das Modell bildet diese Wechselwirkungen nur begrenzt ab und erlaubt 
daher keine eindeutige Aussage darüber, welcher einzelne Faktor für den Erfolg oder Misserfolg eines ERP-Projekts 
ausschlaggebend ist.

Darüber hinaus können die Anforderungen an die einzelnen Fit-Dimensionen je nach Branche, Geschäftsmodell und Unternehmensgröße 
erheblich variieren. In produzierenden Unternehmen kann beispielsweise der Prozess-Fit eine besonders hohe Bedeutung besitzen, 
da Produktionsabläufe, Materialflüsse und Qualitätsanforderungen häufig eng miteinander verbunden sind. In Dienstleistungsunternehmen 
können dagegen Daten-, Kompetenz- oder Strategie-Dimensionen stärker im Vordergrund stehen. Das Modell berücksichtigt solche 
Unterschiede, indem es die einzelnen Dimensionen als flexibel bewertbare Bereiche versteht. Es enthält jedoch keine 
branchenspezifischen Gewichtungen oder allgemein gültigen Schwellenwerte.

Eine weitere Grenze betrifft schwer messbare organisatorische und kulturelle Faktoren. Unternehmenskultur, Veränderungsbereitschaft, 
informelle Machtstrukturen oder Widerstände gegenüber neuen Systemen können den Verlauf einer ERP-Einführung erheblich beeinflussen. 
Diese Faktoren lassen sich jedoch nur eingeschränkt in standardisierte Bewertungskriterien überführen. Ein Unternehmen kann 
beispielsweise formell über ausreichende Kompetenzen und Ressourcen verfügen, während mangelnde Akzeptanz oder Unsicherheit 
bei Mitarbeitenden die tatsächliche Nutzung des ERP-Systems dennoch beeinträchtigen. Solche Aspekte können durch das Modell 
sichtbar gemacht, aber nicht vollständig erfasst oder prognostiziert werden.

Zudem ist ERP lediglich als technologische Grundlage digitaler Transformation zu verstehen. Auch bei einem hohen Fit in allen 
fünf Dimensionen entsteht digitale Transformation nicht automatisch. Sie erfordert zusätzlich eine langfristige strategische 
Ausrichtung, aktive Führung, kontinuierliche Weiterentwicklung von Prozessen sowie die Bereitschaft, bestehende Arbeitsweisen 
zu verändern. Das Modell kann somit kritische Voraussetzungen für eine wirksame ERP-Nutzung strukturieren, ersetzt jedoch keine 
umfassende Digitalisierungsstrategie.

Schließlich liegt eine methodische Limitation der Arbeit darin, dass das ERP-KMU-Fit-Modell literaturbasiert entwickelt wurde. 
Eine empirische Überprüfung anhand realer ERP-Einführungen in unterschiedlichen Branchen und Unternehmen wurde im Rahmen dieser 
Arbeit nicht durchgeführt. Zukünftige Forschung könnte das Modell daher durch Fallstudien, Experteninterviews oder quantitative 
Untersuchungen weiterentwickeln und prüfen, ob einzelne Dimensionen je nach Kontext unterschiedlich gewichtet werden sollten.
